\startfooter{Specialized Project: HK261-DAGD1-041\\
	Market Reconstruction
}
\part{Market Reconstruction} \label{part:market-reconstruction}
Survivorship bias is a well-established problem in empirical financial research: restricting historical analysis to securities that survive until the end of the sample can produce distorted estimates and misleading evidence of return predictability. Brown et al.~\cite{Brown1992} demonstrate that survivorship-truncated samples can materially affect empirical conclusions.

\begin{figure}[H]
	\centering
	\includegraphics*[width=0.6\textwidth]{images/survivorship_bias.png}
	\caption{Abraham Wald's Survivorship Bias Plane}
\end{figure}

For this study, avoiding survivorship bias is complicated by the limited access provided by free-tier market-data services. No single freely accessible source is assumed to provide both a sufficiently long historical security universe and the complete security-level price data required by the study. Consequently, the dataset must be reconstructed by combining complementary sources rather than directly downloaded as a finished historical panel.

\medskip

The pipeline therefore uses the \acf{NYSE} trading calendar as its temporal backbone, retrieves daily historical active-ticker snapshots from \href{https://wrds-www.wharton.upenn.edu/}{Wharton Research Data Services \inlinelogo{images/wrds-logo.png}}, aggregates these observations with Polars to derive ticker availability intervals, and uses those intervals to obtain daily \acs{OHLCV} data from Yahoo Finance. This produces a persistent raw market dataset that preserves historical security availability while remaining feasible under the constraints of the available free-tier \acs{API}s. The resulting data then serve as the foundation for the subsequent integrity checks, point-in-time universe construction, data-quality procedures, and machine-learning pipeline.

\section{U.S. Equity Market Proxy Selection} \label{sec:market-selection}
To construct a representative, investable, and unconfounded U.S. equity universe, this study selects the FTSE Russell 1000 Index as the primary market proxy. Comprising the 1,000 largest U.S.-domiciled securities ranked by total market capitalization within the Russell 3000E universe, the Russell 1000 captures approximately 90\% to 95\% of total investable U.S. equity capitalization while maintaining strict rules-based transparency \cite{ftserussell2026}. At any time $t$, constituents are weighted in proportion to their free-float-adjusted market capitalization $M_{i,t}^{\mathrm{float}} = P_{i,t}(S_{i,t}^{\mathrm{total}} - S_{i,t}^{\mathrm{restricted}})$, yielding the index weight and effective diversification metric:
\begin{equation}
	w_{i,t} = \dfrac{M_{i,t}^{\mathrm{float}}}{\displaystyle\sum_{j=1}^{1000} M_{j,t}^{\mathrm{float}}},
	\qquad
	N_{\mathrm{eff},t} = \dfrac{1}{\displaystyle\sum_{i=1}^{1000} w_{i,t}^2}.
	\label{eq:weights_neff}
\end{equation}

While the top 200 mega-cap constituents account for roughly 78\% of index capitalization (compressing $N_{\mathrm{eff},t}$ and driving the primary market mode), the remaining 800 mid-cap constituents (representing $\approx 22\%$ of index weight, with a median capitalization of \$16.7 billion against a weighted average of \$1.56 trillion as of Q3 2026) provide substantial cross-sectional breadth \cite{ftserussell2026}. For cross-sectional ranking models, expanding the universe from 500 to 1,000 names increases the number of unique pairwise comparisons $\displaystyle\binom{N}{2} = \dfrac{N(N-1)}{2}$ fourfold, from $124{,}750$ to $499{,}500$ per period.

\medskip

The empirical asset-pricing and market-microstructure literature provides three primary justifications for preferring the Russell 1000 over alternative benchmarks such as the S\&P 500 or Russell 3000:

\subsubsection*{Elimination of Committee-Selection Bias}
Unlike the S\&P 500 -- which relies on discretionary committee selection and enforces positive GAAP earnings screens -- Russell 1000 membership is determined mechanically by market-capitalization rank on the prescribed rank day \cite{chang2015, bendavid2019}. As emphasized by \cite{chang2015}, discretionary S\&P 500 additions correlate endogenously with firm profitability, visibility, and expected future performance, confounding cross-sectional return inferences. By contrast, the deterministic Russell 1000/2000 threshold at rank $1{,}000$ creates exogenous variation in benchmark weights. Exploiting this cutoff in a fuzzy regression-discontinuity design (RDD), \cite{chang2015} document that stocks migrating just above rank $1{,}000$ into the top of the value-weighted Russell 2000 experience an approximate $7.3\%$ jump in passive institutional demand and a $5.0\%$ abnormal June return ($t = 2.65$) relative to near-identical stocks remaining at the bottom of the Russell 1000.

\subsubsection*{Well-Documented Reconstitution and Ownership Dynamics}
Because Russell reconstitution follows a predictable schedule -- historically annual in June, and transitioning to a semiannual June/December cycle in 2026 \cite{ftserussell2026} -- its liquidity and ownership shocks are thoroughly parameterized in the literature. Early event studies by \cite{madhavan2003}, \cite{biktimirov2004}, \cite{chen2006}, and \cite{cai2008} show that index rebalancing induces both temporary order-imbalance price pressure (e.g., a $10.78\%$ June addition-minus-deletion spread followed by a $-6.04\%$ July reversal in \cite{madhavan2003}) and permanent shifts in bid-ask spreads and institutional monitoring. Furthermore, \cite{petajisto2011} quantifies the hidden rebalancing drag borne by passive funds during crowded index revisions, while \cite{cremers2009} establish Russell indexes as foundational benchmarks for evaluating active management via Active Share ($\mathrm{AS}_t = \dfrac{1}{2}\displaystyle\sum_i |w_{i,t}^{\mathrm{fund}} - w_{i,t}|$). More recently, researchers have used Russell's mechanical assignment to identify how passive and ETF ownership improve corporate governance \cite{appel2016}, propagate return comovement \cite{chen2016, da2018}, and elevate daily volatility and short-horizon negative autocorrelation -- yielding a risk premium of up to 56 basis points per month for high-ETF-ownership stocks \cite{bendavid2018}.

\subsubsection*{Institutional Liquidity and Arbitrage Efficiency vs.\ the Russell 3000}
While the broader Russell 3000 captures roughly 98\% of the U.S. market by adding the 2,000 small-cap constituents of the Russell 2000 \cite{ftserussell2026}, the bottom 2,000 names contribute only 3\% to 8\% of market capitalization while imposing severe implementation friction. Using Russell reconstitution as an instrument, \cite{patatoukas2022} demonstrate that large- and mid-cap Russell 1000 securities exhibit low arbitrage constraints and unimpaired price discovery, whereas small- and micro-cap stocks suffer from binding short-sale and liquidity frictions. Under a standard square-root market-impact cost function for trade size $|Q_{i,t}|$ relative to average daily dollar volume $\mathrm{ADV}_{i,t}$ and volatility $\sigma_{i,t}$:
\begin{equation}
	MI_{i,t} = \eta_i \sigma_{i,t} \left( \dfrac{|Q_{i,t}|}{\mathrm{ADV}_{i,t}} \right)^{1/2},
	\label{eq:market_impact}
\end{equation}
truncating the universe at rank $1{,}000$ avoids the steep market-impact decay, wide spreads, and elevated delisting bias of the small-cap tail while preserving the economically dominant core of the U.S. equity market.

\section{Data Acquisition and Processing Pipeline} \label{sec:data-acquisition}
To construct a survivorship-bias-free dataset, we implement a three-stage data acquisition and normalization pipeline. Because standard financial databases frequently overwrite historical index memberships with current constituents, reconstructing point-in-time constituent lists is essential for empirical validity. Our workflow systematically captures annual snapshots of the Russell 1000 Index following its late-June reconstitution and links these constituents to fundamental and market records.

\subsection{Historical Index Constituent Extraction}
The Russell 1000 Index undergoes an annual reconstitution on the final Friday of June, establishing the definitive constituent list for the subsequent twelve months. To capture true point-in-time index membership across our 2000--2026 sample period, we retrieve archival constituent snapshots directly from the \href{web.archive.org}{Internet Archive's Wayback Machine} immediately following each June reconstitution phase.

\medskip

Because institutional data hosting and corporate ownership have migrated across platforms over the past two decades, we query multiple archival domains to ensure complete longitudinal coverage:
\begin{itemize}
	\item \textbf{Official Index Providers:} Primary constituent lists and reconstitution reports published across current and legacy provider portals, including \href{https://lseg.com/}{London Stock Exchange Group}, \href{https://ftserussell.com/}{FTSE Russell}, \href{https://russellinvestments.com/}{Russell Investments} and \href{https://russell.com/}{russell.com}, and the defunct early-2000s \href{https://www.russellmellon.com/}{Russell/Mellon} analytical services domain.
	\item \textbf{Institutional ETF Holdings:} Post-reconstitution holdings disclosures from \href{https://ishares.com/}{BlackRock} for the iShares Russell 1000 ETF (IWB), which serve as a high-fidelity proxy and cross-validation source when direct provider tables are unavailable.
\end{itemize}
Due to changes in web standards over the 2000--2026 sample horizon, the 27 extracted raw files span heterogeneous file architectures, including unstructured Portable Document Format (\texttt{.pdf}) fact sheets, Microsoft Excel (\texttt{.xls}) and comma-separated values (\texttt{.csv}) spreadsheets, and JavaScript Object Notation (\texttt{.json}) API payloads (summarized in Table~\ref{tab:russell_raw_sources}). We implement a custom parsing pipeline to strip document metadata, headers, and footnotes, standardizing each annual extraction into a clean, line-separated plain-text file containing company names and primary ticker symbols for each reconstitution year $t$.
\begin{table}[H]
	\small
	\centering
	\caption{Raw Russell 1000 Constituent Snapshots}
	\label{tab:russell_raw_sources}
	\begin{tabular}{lcll}
		\toprule
		\textbf{File Format} & \textbf{Count} & \textbf{Reconstitution Cohort Years} & \textbf{Primary Source Structure}    \\
		\midrule
		PDF Document         & 21             & 2000--2007, 2009--2018, 2020--2022   & Multi-page provider fact sheets      \\
		JSON Payload         & 3              & 2019, 2023, 2024                     & Structured web API responses         \\
		Excel Workbook       & 2              & 2025, 2026                           & Tabular holdings spreadsheets        \\
		CSV File             & 1              & 2008                                 & Delimited constituent table          \\
		\midrule
		\textbf{Total}       & \textbf{27}    & \textbf{2000--2026 (Complete Panel)} & \textbf{Normalized to \texttt{.txt}} \\
		\bottomrule
	\end{tabular}
\end{table}

\subsection{WRDS Query Execution and Historical Data Retrieval} \label{subsec:wrds-query}
With the standardized annual constituent lists established, the second phase of the pipeline extracts longitudinal stock records via the \href{https://wrds-www.wharton.upenn.edu/}{Wharton Research Data Services \inlinelogo{images/wrds-logo.png}} web query interface, specifically targeting the Center for Research in Security Prices (CRSP) Stock Database (\texttt{crsp\_a\_stock}). To preserve the annual reconstitution structure and prevent look-ahead or survivorship bias, we execute separate cohort-specific queries for each reconstitution cycle rather than pooling tickers across the entire sample period.

\begin{figure}[H]
	\centering
	\includegraphics*[width=0.9\textwidth]{images/wrds-flow.png}
	\caption{Historical Constituent Data Pipeline.}
\end{figure}

For each reconstitution year $t$, the WRDS web query execution follows a standardized four-step protocol:
\begin{enumerate}[topsep=0pt, parsep=0pt]
	\item[Step 1] \textbf{Date Range:} Select the 12-month active cycle following reconstitution, from late June of year $t$ to late June of year $t+1$ (e.g., \texttt{2023-06-30} to \texttt{2024-06-30}).
	\item[Step 2] \textbf{Code Upload:} Set the format to ticker symbols and upload the standardized plain-text file (\texttt{.txt}, one ticker per line) using the WRDS interface.
	\item[Step 3] \textbf{Variables:} Select 61 Identifying and Time Series variables to capture market data and security-level metadata.
	\item[Step 4] \textbf{Output:} Export the query results as comma-delimited text (\texttt{.csv}) in a compressed \texttt{.zip} archive using \texttt{YYYY-MM-DD} dates, then download directly from the WRDS portal.
\end{enumerate}

\section{Data Integrity}